\section{Validación del modelo conceptual}

Antes de su implementación, el modelo conceptual fue contrastado con el
relevamiento realizado y con la estructura física de la guardia.

\subsection{Correspondencia con el sistema real}

El flujo descripto en la sección~2.4 fue verificado contra el circuito que
sigue efectivamente un paciente en la guardia: ingreso por la puerta
principal, admisión administrativa, espera en la sala común, llamado a
triage, retorno a la sala, llamado a consultorio y egreso. La derivación
directa de los pacientes de ambulancia al Shock Room y la identificación de
casos críticos en el acceso también fueron confirmadas durante el
relevamiento.

\subsection{Decisiones de abstracción y su justificación}

\begin{itemize}
  \item \textbf{Espera en sala común.} Las dos colas del sistema --- triage y
  consultorio --- comparten un mismo espacio físico. El modelo conceptual las
  representa como colas lógicas separadas sobre una misma localización, lo
  que permite medir el hacinamiento agregado sin perder la distinción entre
  ambas esperas.
  \item \textbf{Prioridad no expulsiva.} Un paciente que ya comenzó su
  atención no es interrumpido por la llegada de otro más grave. Esta decisión
  refleja la práctica del servicio.
  \item \textbf{Internación como punto de salida.} Los procesos internos de
  las salas de internación no influyen sobre las métricas de interés, por lo
  que su exclusión no compromete la validez del modelo respecto del problema
  planteado.
  \item \textbf{Cola ciega previa al triage.} El modelo representa
  explícitamente que la priorización sólo existe aguas abajo de la
  clasificación, lo que resulta esencial para interpretar correctamente el
  escenario alternativo.
\end{itemize}

\subsection{Limitaciones conocidas del modelo conceptual}

\begin{itemize}
  \item No se modela el abandono de pacientes que se retiran sin ser
  atendidos. De existir en el sistema real, el modelo sobreestima la carga.
  \item No se modela la reclasificación de pacientes que empeoran durante la
  espera.
  \item La tasa de error de clasificación es independiente de la carga de
  trabajo.
  \item La exclusión de traumatología, radiografías y medicación concentra en
  cuatro consultorios una demanda que en el sistema real se reparte entre más
  sectores.
\end{itemize}

